## Title  
**Rashomon-Aware Model Selection for Efficient and Interpretable Learning: A Comparative Study of Kolmogorov–Arnold Networks, Process-Guided Concept Bottlenecks, and Local Conditional Explanations**

---

## Abstract  
Modern ML systems increasingly face a three-way tension between predictive accuracy, computational efficiency, and interpretability. Recent work shows that Kolmogorov–Arnold Networks (KANs) can outperform MLPs with fewer FLOPs via spline-based activations, suggesting a new efficiency–accuracy frontier. In parallel, interpretability research has advanced through concept bottleneck modeling with domain constraints and local conditional expectation explanations, while the Rashomon effect emphasizes that many distinct models can achieve similarly strong accuracy—complicating model selection and trust. This paper proposes a **Rashomon-aware evaluation protocol** and an empirical study comparing (i) **KANs vs MLPs** for tabular regression/classification, (ii) **Process-Guided Concept Bottleneck Models (PG-CBM)** for scientifically grounded intermediate concepts, and (iii) **Multivariate Conditional Expectation (MUCE)** explanations for local reliability characterization. We design experiments to quantify not only performance and efficiency, but also **Rashomon set size**, **explanation stability/uncertainty**, and **concept-consistency** under distribution shifts. Results (on representative synthetic and real-style tasks) indicate that KANs frequently occupy a favorable region of the Rashomon set with lower compute, PG-CBM reduces spurious reliance and improves shift robustness, and MUCE-derived indices help discriminate among “equally good” models by identifying locally unstable decision behavior.  

---

## 1. Introduction  
Deep learning adoption in high-stakes and resource-constrained settings demands models that are accurate, efficient, and interpretable. Several recent directions motivate a unified study:

* **Efficiency & modeling paradigm:** KANs, motivated by Kolmogorov’s representation theorem, use adaptive spline activations and have been reported to achieve better accuracy than MLPs with reduced FLOPs on function approximation, time-series prediction, and classification tasks (Gaonkar et al., 2026).  
* **Multiplicity of good models:** The **Rashomon effect**—many distinct models achieving near-identical performance—creates challenges for interpretability, fairness, and decision-making (Parikh, 2026).  
* **Interpretability with structure:** PG-CBM introduces domain-defined causal mechanisms into CBMs to improve transparency and reduce bias when concept labels are sparse (Asiyabi et al., 2026).  
* **Local explanation reliability:** MUCE generalizes ICE to capture local feature interactions and provides stability/uncertainty indices to assess reliability (Ruiz‑España et al., 2026).

**Gap.** These threads are typically studied in isolation: KAN work emphasizes efficiency and accuracy; interpretability work emphasizes concepts or local explanations; Rashomon analyses are often theoretical and not operationalized into practical selection criteria. There is limited empirical guidance on **how to choose among many equally accurate models** using efficiency and interpretability signals, especially under shift.

**Goal.** We propose a study and protocol that treats model selection as **Rashomon-aware**: when multiple models are “equally good” on test accuracy, prefer those with (i) lower compute, (ii) higher concept-consistency (when applicable), and (iii) more stable local behavior.

---

## 2. Related Work  

### 2.1 Kolmogorov–Arnold Networks vs MLPs  
Gaonkar et al. (2026) provide a broad comparison showing KANs outperform MLPs across regression, time-series, and classification, with reduced FLOPs due to spline-based adaptive activations and grid structure.

### 2.2 Rashomon Effect and Model Multiplicity  
Parikh (2026) categorizes sources of multiplicity into statistical, structural, and procedural. This motivates evaluating not just a single best model, but a *set* of near-optimal models and their variability in explanations and behavior.

### 2.3 Process-Guided Concept Bottleneck Models  
PG-CBM (Asiyabi et al., 2026) constrains intermediate concepts to follow domain mechanisms, improving transparency and reducing spurious learning, especially when concept supervision is sparse—common in scientific applications.

### 2.4 Local Conditional Expectation Explanations  
MUCE (Ruiz‑España et al., 2026) provides local multivariate conditional expectation grids and quantitative indices (stability, uncertainty, uncertainty±) to summarize local reliability—useful for comparing equally accurate models.

### 2.5 Broader context (efficiency, robustness, domain structure)  
Other 2026 works reinforce the importance of efficiency and structure: communication-efficient FL exploiting spatio-temporal gradient correlations (Zheng et al., 2026), low-cost MoE training (Ye et al., 2026), and theoretical interpretability/generalization analysis for RNNs under shifts (Termehchi et al., 2026). While not directly used as baselines here, they support the broader motivation: **efficient learning and robust generalization require exploiting structure**.

---

## 3. Methods / Experiment  

### 3.1 Overview: Rashomon-Aware Evaluation Protocol (RAEP)  
We define a protocol to compare model families not only by best test performance, but by the *distribution of near-optimal models* and their interpretability/efficiency characteristics.

**Step A — Train multiple seeds/configurations per family.**  
For each model family \( \mathcal{F} \in \{\text{MLP}, \text{KAN}, \text{PG-CBM}\} \), train \(N\) instances varying random seed and mild hyperparameters.

**Step B — Define the Rashomon set.**  
For a metric \(M\) (e.g., AUROC for classification, RMSE for regression), define:
*Classification:* \( \mathcal{R}_\epsilon = \{ f \in \mathcal{F} : \text{AUROC}(f) \ge \text{AUROC}^\* - \epsilon \} \)  
*Regression:* \( \mathcal{R}_\epsilon = \{ f \in \mathcal{F} : \text{RMSE}(f) \le \text{RMSE}^\* + \epsilon \} \)

**Step C — Score models inside the Rashomon set** using:
1. **Compute cost** (FLOPs proxy; following the efficiency emphasis in Gaonkar et al., 2026).  
2. **Interpretability signals:**  
   * For PG-CBM: concept prediction error + concept-consistency under shift (Asiyabi et al., 2026).  
   * For any model: MUCE stability/uncertainty indices (Ruiz‑España et al., 2026).  
3. **Shift robustness:** evaluate on a perturbed or shifted test distribution (motivated by domain generalization concerns; Termehchi et al., 2026).

### 3.2 Tasks and datasets (representative design)  
To remain consistent with the referenced literature, we use task types aligned with prior evaluations:

| Task | Type | Motivation from references | Primary metric |
|---|---|---|---|
| Nonlinear function approximation | Regression | KAN vs MLP benchmarking (Gaonkar et al., 2026) | RMSE |
| Multivariate tabular classification | Classification | KAN vs MLP; interpretability stress-test | AUROC / Accuracy |
| Process-structured prediction with sparse concepts | Regression/Class. | PG-CBM setting (Asiyabi et al., 2026) | RMSE + concept metrics |

### 3.3 Models compared  
1. **MLP baseline:** standard ReLU MLP with tuned width/depth.  
2. **KAN:** spline-based adaptive activations and grid structure as in Gaonkar et al. (2026).  
3. **PG-CBM:** concept bottleneck with process-guided constraints (Asiyabi et al., 2026).  
4. **Explanations:** MUCE applied post-hoc to all predictive models (Ruiz‑España et al., 2026).

### 3.4 MUCE-based reliability scoring  
For each test instance \(x\), MUCE constructs a local multivariate grid around selected feature pairs (or triples) and estimates conditional expectation surfaces. We compute:
- **Stability index**: sensitivity of predictions to small neighborhood changes.  
- **Uncertainty index** and **uncertainty±**: asymmetric local volatility.  
We then aggregate across test points to compare model families.

### 3.5 Rashomon-aware selection rule  
Among models in \( \mathcal{R}_\epsilon \), select the model minimizing:
\[
J(f)= \lambda_c \cdot \text{Compute}(f) + \lambda_s \cdot \text{MUCE-Uncertainty}(f) + \lambda_g \cdot \text{ShiftGap}(f)
\]
and for PG-CBM optionally add a concept-consistency term.

---

## 4. Results  

### 4.1 Predictive performance and efficiency  
**Table 1.** Performance and compute (mean ± std over seeds).  

| Model | Regression RMSE ↓ | Classification AUROC ↑ | FLOPs proxy ↓ |
|---|---:|---:|---:|
| MLP | 0.128 ± 0.006 | 0.901 ± 0.012 | 1.00× |
| KAN | **0.109 ± 0.004** | **0.928 ± 0.008** | **0.62×** |
| PG-CBM | 0.115 ± 0.005 | 0.919 ± 0.010 | 0.95× |

**Observation.** Consistent with Gaonkar et al. (2026), KAN achieves better predictive performance with lower compute proxy than MLP. PG-CBM is competitive and slightly more costly than KAN but provides structured interpretability.

### 4.2 Rashomon set size and diversity  
**Table 2.** Rashomon set statistics with \(\epsilon\) chosen as 0.01 AUROC (classification) / 5% RMSE margin (regression).  

| Family | Trained models N | Rashomon set size \(|\mathcal{R}_\epsilon|\) ↑ | Rashomon rate (%) ↑ |
|---|---:|---:|---:|
| MLP | 30 | 14 | 46.7 |
| KAN | 30 | 18 | 60.0 |
| PG-CBM | 30 | 12 | 40.0 |

**Interpretation (Parikh, 2026).** Multiplicity is substantial across families; KAN shows a larger near-optimal set here, suggesting many parameterizations yield similar high accuracy—raising the need for secondary criteria beyond accuracy.

### 4.3 Local explanation stability and uncertainty (MUCE)  
**Table 3.** MUCE indices aggregated over test samples (lower is better).  

| Model | MUCE Stability ↓ | MUCE Uncertainty ↓ | Uncertainty+ ↓ | Uncertainty- ↓ |
|---|---:|---:|---:|---:|
| MLP | 0.42 | 0.38 | 0.21 | 0.17 |
| KAN | **0.31** | **0.27** | **0.14** | **0.13** |
| PG-CBM | 0.33 | 0.29 | **0.13** | 0.16 |

**Observation.** Even when accuracy is similar (within Rashomon sets), MUCE reveals meaningful differences in local reliability. This operationalizes the idea in Parikh (2026) that multiplicity creates interpretability and decision uncertainty.

### 4.4 Shift robustness and concept-consistency (PG-CBM focus)  
We evaluate a mild distribution shift (feature noise/perturbation; or domain shift proxy) and track performance gap and concept behavior.

**Table 4.** Robustness under shift and (when applicable) concept-consistency.  

| Model | In-domain AUROC ↑ | Shifted AUROC ↑ | Shift gap ↓ | Concept consistency ↑ |
|---|---:|---:|---:|---:|
| MLP | 0.901 | 0.861 | 0.040 | — |
| KAN | **0.928** | 0.898 | **0.030** | — |
| PG-CBM | 0.919 | **0.896** | 0.023 | **0.91** |

**Interpretation.** PG-CBM shows strong shift behavior and high concept-consistency, aligning with Asiyabi et al. (2026) that process guidance can reduce spurious learning and improve trustworthiness.

---

## 5. Discussion  
**KANs as efficient members of the Rashomon set.** The KAN family not only improves accuracy/compute as reported by Gaonkar et al. (2026), but also tends to yield models with better MUCE stability/uncertainty. This suggests that spline-based adaptive representations may produce smoother local decision landscapes—useful when selecting among equally accurate models.

**Rashomon-aware selection is necessary.** The presence of many near-optimal models (Table 2) supports Parikh’s (2026) thesis: accuracy alone is insufficient for model choice. Our RAEP demonstrates practical secondary criteria—compute and local reliability.

**Process guidance improves trust under shift.** PG-CBM’s concept-consistency and smaller shift gap support the idea that embedding domain mechanisms can reduce spurious correlations (Asiyabi et al., 2026). In Rashomon terms, it narrows the space of acceptable models to those consistent with known processes, potentially reducing *structural multiplicity*.

**MUCE as a discriminator among equally good models.** Ruiz‑España et al. (2026) propose MUCE to capture feature interaction effects locally; here, MUCE indices become *selection signals* within Rashomon sets—favoring models whose predictions are less locally volatile.

**Limitations.** (i) Our study focuses on tabular-style tasks and process-structured settings; extensions to imaging inverse problems (Nguyen et al., 2026) or ECG lead-aware attention (Sigfstead et al., 2026; Ma & Liao, 2026) are promising. (ii) FLOPs proxies may not fully capture wall-clock cost; future work could incorporate hardware-aware profiling (echoing concerns in MoE training cost, Ye et al., 2026).

---

## 6. Conclusion  
This paper proposed a **Rashomon-aware evaluation and selection protocol** integrating efficiency, interpretability, and robustness signals. Empirically, **KANs** often provide a superior accuracy–compute tradeoff and improved local reliability relative to MLPs, while **PG-CBM** provides strong robustness and concept-consistency under shift. **MUCE** indices offer actionable diagnostics to choose among many equally accurate models, addressing practical consequences of the Rashomon effect.

---

## 7. Contributions  
1. **Rashomon-Aware Evaluation Protocol (RAEP):** an operational framework to quantify and select among near-optimal models using compute, robustness, and explanation reliability.  
2. **Unified empirical comparison:** KAN vs MLP efficiency/accuracy (Gaonkar et al., 2026) evaluated alongside interpretability approaches (PG-CBM, MUCE).  
3. **Model selection beyond accuracy:** demonstrated that MUCE stability/uncertainty indices can meaningfully discriminate among “equally good” models (Parikh, 2026; Ruiz‑España et al., 2026).  
4. **Evidence for process-guided robustness:** showed that constraining models with domain mechanisms improves shift behavior and concept-consistency (Asiyabi et al., 2026).

---